\documentclass[10pt,a4paper]{amsart}
\usepackage[margin=19mm]{geometry}
\usepackage{amsmath,amssymb,mathtools}
\usepackage[scr=rsfs,cal=euler]{mathalfa}
\usepackage{xfrac}
\usepackage[unicode,hidelinks,bookmarks=false]{hyperref}
\newcommand{\ie}{i.e.\ }
\newcommand{\cf}{cf.\ }
\newcommand{\resp}{respectively\ }
\newcommand{\Subsection}{\subsection}
\providecommand{\nobreakdash}{\nobreak\hbox{-}}
\setlength{\emergencystretch}{2em}
\multlinegap0pt
\usepackage{bm,mathtools}
\usepackage{booktabs}
\usepackage{multirow}
\usepackage{algorithm}
\usepackage{algpseudocode}
\usepackage{color}
\usepackage{enumitem}
\usepackage{multicol}
\newtheorem{theorem}{Theorem}
\newtheorem{remark}{Remark}
\newtheorem{definition}{Definition}
\newtheorem{proposition}{Proposition}
\newcommand{\R}{\mathbb{R}}
\newcommand{\norm}[1]{\left\|#1\right\|}
\newcommand{\grad}{\nabla}
\newcommand{\argmin}{\operatorname*{arg\,min}}
\newcommand{\Lop}{\mathcal{L}}
\begin{document}
\title[IGA-ODIL: Optimizing DIscretre robust Loss with Isogeometric]{IGA-ODIL: Optimizing DIscretre robust Loss with Isogeometric Analysis to solve forward and inverse problems faster using machine learning tools}
\author{Maciej Paszyński, Tomasz S{ł}użalec}
\date{}
\maketitle
\noindent\emph{Source and licence.} IGA-ODIL: Optimizing DIscretre robust Loss with Isogeometric Analysis to solve forward and inverse problems faster using machine learning tools. Version arXiv:2605.30272v1 (CC BY 4.0). This material is distributed under the Creative Commons Attribution 4.0 International licence, http://creativecommons.org/licenses/by/4.0/, as recorded in the arXiv OAI licence metadata for this exact version and shown on the abstract page https://arxiv.org/abs/2605.30272v1. Full rights notice: licenses/arxiv-2605.30272-CC-BY-4.0.txt.\par\smallskip
\noindent\textit{Editorial scope.} Counted unit: the original \texttt{theorem} environment \texttt{thm:robust\_iga\_odil} at source lines 1724--1896 together with its complete proof at lines 1898--2114; the delivered document keeps source lines 1722--2115 so that every referenced label and every citation resolves inside it. Native editable source is the paper's own concatenated TeX; the AMS wrapper is editorial formatting only. No publisher class, upstream script or unrelated artwork is redistributed.
\section{Convergence of IGA-ODIL}


\begin{theorem}[Convergence of robust spline-based residual minimization]
\label{thm:robust_iga_odil}

Let $\Omega \subset \mathbb{R}^d$ be a bounded Lipschitz domain, and consider the boundary value problem
\begin{equation}
\mathcal{L}u = f
\quad \text{in } \Omega,
\label{eq:general_pde}
\end{equation}
supplemented with boundary conditions such that the problem is well posed.

Assume that:

\begin{enumerate}

\item
$\mathcal{L}:H^m(\Omega)\to L^2(\Omega)$
is a linear differential operator of order
$
s\le m.
$

\item
The exact solution satisfies
$
u\in H^{p+1}(\Omega).
$

\item
The spline approximation space
$
V_h\subset H^m(\Omega)
$
consists of tensor-product B-splines of degree
$
p
$
constructed on a quasi-uniform mesh with characteristic size
$
h.
$

\item
The collocation points
$
\mathcal{X}_M=\{x_k\}_{k=1}^M\subset\Omega
$
are quasi-uniform with fill distance
$
h_c
:=
\sup_{x\in\Omega}
\min_{x_k\in\mathcal{X}_M}
\|x-x_k\|,
$
satisfying
$
h_c\lesssim h.
$

\item
The discrete robust residual minimization problem
\begin{equation}
u_h
=
\arg\min_{v_h\in V_h}
\Phi_h(v_h)
\label{eq:robust_discrete_problem}
\end{equation}
is defined by the functional
\begin{equation}
\Phi_h(v_h)
=
\frac12
R(v_h)^T
G^{-1}
R(v_h),
\label{eq:robust_functional}
\end{equation}
where
\begin{equation}
R(v_h)
=
\left[
\mathcal{L}v_h(x_1)-f(x_1),
\dots,
\mathcal{L}v_h(x_M)-f(x_M)
\right]^T
\in\mathbb{R}^M,
\end{equation}
and
$
G\in\mathbb{R}^{M\times M}
$
is a symmetric positive definite Gram matrix.

\item
The operator $\mathcal{L}$ satisfies the stability estimate
\begin{equation}
\|w\|_{H^m(\Omega)}
\le
C_{\mathrm{stab}}
\|\mathcal{L}w\|_{L^2(\Omega)}
\quad
\forall w\in H_0^m(\Omega).
\label{eq:stability_assumption_theorem}
\end{equation}

\item
The collocation sampling satisfies the norm equivalence
\begin{equation}
\|v_h\|_{L^2(\Omega)}^2
\sim
\sum_{k=1}^{M}
|v_h(x_k)|^2
\quad
\forall v_h\in V_h,
\label{eq:sampling_equivalence_theorem}
\end{equation}
with constants independent of $h$.

\item
The Gram matrix induces a uniformly equivalent discrete norm:
\begin{equation}
c_G\|r\|_2^2
\le
r^TG^{-1}r
\le
C_G\|r\|_2^2
\quad
\forall r\in\mathbb{R}^M,
\label{eq:gram_equivalence_theorem}
\end{equation}
with constants independent of $h$.

\end{enumerate}

Then the robust IGA--ODIL approximation satisfies the error estimate
\begin{equation}
\|u-u_h\|_{H^m(\Omega)}
\le
C
\|R(u_h)\|_G,
\label{eq:error_equivalence}
\end{equation}
where
\begin{equation}
\|R(u_h)\|_G
:=
\left(
R(u_h)^TG^{-1}R(u_h)
\right)^{1/2}.
\label{eq:robust_norm}
\end{equation}

Moreover,
\begin{equation}
\|R(u_h)\|_G
\le
C
h^{p+1-s}
|u|_{H^{p+1}(\Omega)},
\label{eq:residual_convergence_theorem}
\end{equation}
and therefore
\begin{equation}
\|u-u_h\|_{H^m(\Omega)}
=
\mathcal{O}(h^{p+1-s}).
\label{eq:final_rate_theorem}
\end{equation}

\end{theorem}

\begin{proof}

\noindent
\textbf{Step 1: Spline approximation property.}

Since
$
u\in H^{p+1}(\Omega),
$
classical spline approximation theory
(cf.\ Ciarlet~\cite{ciarlet2002finite},
Schumaker~\cite{schumaker2007spline},
Bazilevs et al.~\cite{bazilevs2006isogeometric})
implies the existence of a spline interpolant
$
\Pi_hu\in V_h
$
such that
\begin{equation}
\|u-\Pi_hu\|_{H^m(\Omega)}
\le
Ch^{p+1-m}
|u|_{H^{p+1}(\Omega)}.
\label{eq:proof_interp}
\end{equation}

Since $\mathcal{L}$ is a differential operator of order $s$,
standard continuity estimates yield
\begin{equation}
\|\mathcal{L}(u-\Pi_hu)\|_{L^2(\Omega)}
\le
Ch^{p+1-s}
|u|_{H^{p+1}(\Omega)}.
\label{eq:proof_operator_error}
\end{equation}


\noindent
\textbf{Step 2: Minimality of the robust residual functional.}

Since
$
u_h
=
\arg\min_{v_h\in V_h}
\Phi_h(v_h),
$
we have
\begin{equation}
\Phi_h(u_h)
\le
\Phi_h(\Pi_hu).
\label{eq:proof_minimality}
\end{equation}

Using the definition of the robust functional,
$
\Phi_h(v_h)
=
\frac12
R(v_h)^TG^{-1}R(v_h),
$
this becomes
\begin{equation}
R(u_h)^TG^{-1}R(u_h)
\le
R(\Pi_hu)^TG^{-1}R(\Pi_hu).
\label{eq:proof_robust_minimality}
\end{equation}

Since the exact solution satisfies
$
\mathcal{L}u=f,
$
the residual associated with the interpolant satisfies
$
R(\Pi_hu)
=
\left[
\mathcal{L}(\Pi_hu-u)(x_k)
\right]_{k=1}^{M}.
$

Using the Gram norm equivalence
\eqref{eq:gram_equivalence_theorem},
\begin{equation}
R(\Pi_hu)^TG^{-1}R(\Pi_hu)
\le
C_G
\sum_{k=1}^{M}
|\mathcal{L}(\Pi_hu-u)(x_k)|^2.
\label{eq:proof_gram_upper}
\end{equation}

Applying the sampling equivalence
\eqref{eq:sampling_equivalence_theorem},
\begin{equation}
\sum_{k=1}^{M}
|\mathcal{L}(\Pi_hu-u)(x_k)|^2
\le
C
\|\mathcal{L}(\Pi_hu-u)\|_{L^2(\Omega)}^2.
\label{eq:proof_sampling}
\end{equation}

Combining
\eqref{eq:proof_operator_error},
\eqref{eq:proof_gram_upper},
and
\eqref{eq:proof_sampling}
yields
\begin{equation}
R(u_h)^TG^{-1}R(u_h)
\le
Ch^{2(p+1-s)}
|u|_{H^{p+1}(\Omega)}^2.
\label{eq:proof_residual_bound_sq}
\end{equation}

Taking square roots gives
\begin{equation}
\|R(u_h)\|_G
\le
Ch^{p+1-s}
|u|_{H^{p+1}(\Omega)}.
\label{eq:proof_residual_bound}
\end{equation}

This proves
\eqref{eq:residual_convergence_theorem}.


\noindent
\textbf{Step 3: Stability estimate.}

Define the approximation error
$
e_h:=u-u_h.
$
Using linearity of the operator,
$
\mathcal{L}e_h
=
f-\mathcal{L}u_h.
$
Applying the stability estimate
\eqref{eq:stability_assumption_theorem},
\begin{equation}
\|e_h\|_{H^m(\Omega)}
\le
C_{\mathrm{stab}}
\|f-\mathcal{L}u_h\|_{L^2(\Omega)}.
\label{eq:proof_stability}
\end{equation}

Using the sampling equivalence,
\begin{equation}
\|f-\mathcal{L}u_h\|_{L^2(\Omega)}^2
\sim
\sum_{k=1}^{M}
|\mathcal{L}u_h(x_k)-f(x_k)|^2.
\label{eq:proof_sampling_residual}
\end{equation}

Applying the Gram equivalence
\eqref{eq:gram_equivalence_theorem},
\begin{equation}
\sum_{k=1}^{M}
|\mathcal{L}u_h(x_k)-f(x_k)|^2
\le
c_G^{-1}
R(u_h)^TG^{-1}R(u_h).
\label{eq:proof_gram_lower}
\end{equation}

Therefore,
\begin{equation}
\|f-\mathcal{L}u_h\|_{L^2(\Omega)}
\le
C
\|R(u_h)\|_G.
\label{eq:proof_residual_norm}
\end{equation}

Substituting into
\eqref{eq:proof_stability}
gives
\begin{equation}
\|u-u_h\|_{H^m(\Omega)}
\le
C
\|R(u_h)\|_G.
\label{eq:proof_error_estimate}
\end{equation}

Finally, combining
\eqref{eq:proof_error_estimate}
with
\eqref{eq:proof_residual_bound}
yields
\begin{equation}
\|u-u_h\|_{H^m(\Omega)}
\le
Ch^{p+1-s}
|u|_{H^{p+1}(\Omega)}.
\end{equation}

Hence,
$
\|u-u_h\|_{H^m(\Omega)}
=
\mathcal{O}(h^{p+1-s}),
$
which proves
\eqref{eq:final_rate_theorem}.

\end{proof}


\begin{thebibliography}{99}
\bibitem[bazilevs2006isogeometric]{bazilevs2006isogeometric} Y.~Bazilevs, L.~Beir\~ao da Veiga, J.~A.~Cottrell, T.~J.~R.~Hughes, G.~Sangalli, Isogeometric analysis: Approximation, stability and error estimates for $h$-refined meshes, \emph{Mathematical Models and Methods in Applied Sciences}, 17 (2006) 1031--1090.
\bibitem[ciarlet2002finite]{ciarlet2002finite} P.~G.~Ciarlet, \emph{The Finite Element Method for Elliptic Problems}, SIAM, Philadelphia, 2002.
\bibitem[schumaker2007spline]{schumaker2007spline} L.~L.~Schumaker, \emph{Spline Functions: Basic Theory}, 3rd ed., Cambridge University Press, Cambridge, 2007.
\end{thebibliography}
\end{document}
